\section{Introduction}
\label{sec:intro}

Fine-tuning a pre-trained transformer with LoRA~\cite{hu2022lora} is
fundamentally different from training from scratch: the model arrives with
strong priors, and the informativeness of each training sample changes as
fine-tuning proceeds.
Some examples become \emph{mastered} --- their loss approaches zero and their
gradient contribution is negligible.
Others are \emph{conflicting} --- labels are irreconcilable with the pre-trained
prior (e.g.\ label noise), so their loss stagnates.
Between these two groups sits a shifting minority at the \emph{adaptation
frontier}: examples where loss is actively declining and gradients are large and
directionally consistent.
Standard SFT with LoRA treats all three groups equally, diluting each mini-batch
gradient with near-zero contributions from mastered samples and noisy
contributions from conflicting ones.

Prior work on this problem leaves a gap.
Curriculum and self-paced learning~\cite{bengio2009curriculum,kumar2010self}
weight samples by loss \emph{magnitude}, which conflates frontier with
irreconcilable noise --- both exhibit high loss, but only the former benefits
from upweighting.
Meta-reweighting~\cite{ren2018learning,shu2019meta} learns optimal weights via
second-order gradients, but requires a clean held-out set and roughly doubles
compute.
Dataset Cartography~\cite{swayamdipta2020dataset} maps the ``ambiguous'' stratum
--- the online analogue of our frontier --- but requires a dedicated
preliminary training pass.

We propose \textbf{EWL (Example Weighting via Learning Progress)}, a method that
identifies the adaptation frontier online with no additional data or compute
beyond the training run itself.
EWL tracks per-sample loss velocity as a relative progress signal via
exponential moving averages (EMA), then modulates it by a \emph{LoRA gradient
proxy} --- the mean Frobenius norm of adapter gradients.
The proxy acts as an adaptive temperature: large during rapid adaptation, it
sharpens the weight distribution toward frontier samples; near convergence, it
decays and gracefully recovers uniform weighting.
This gating is critical: without it, the raw progress signal applies aggressive
reweighting before the adapter has accumulated meaningful history, destabilising
training.
The per-sample weights are applied via a softmax-weighted loss with a single
temperature $\tau$.
EWL requires no meta-gradients, no held-out data, and adds no inference cost.

Experiments on three fine-grained visual benchmarks using ViT-Small/16 with LoRA
reveal two findings.
First, the LoRA gradient proxy is not optional: ablating it causes a $-6.76\%$
accuracy drop on FGVC-Aircraft, demonstrating that ungated reweighting
destabilises training on tasks with densely packed decision boundaries.
Second, EWL's noise robustness is \emph{dataset-conditional}: on Aircraft, where
extreme inter-class similarity sustains a wide adaptation frontier, EWL
improves over SFT by $+4.3$--$6.8\%$ under 10--40\% label noise; on CUB-200 and
Stanford Dogs, where classes are more separable, EWL is harmful under noise
($-0.9\%$ to $-6.6\%$).
On clean data, EWL matches SFT within standard deviation at a $5.9\times$
step-time overhead with zero additional GPU memory.

\noindent\textbf{Contributions:} (i)~\textbf{EWL}, a single-pass, online sample
weighting method for LoRA fine-tuning based on per-sample loss velocity gated
by a LoRA gradient proxy; (ii)~empirical evidence that the gradient proxy is
\textbf{essential for training stability} --- ablating it causes severe
degradation on narrow-boundary tasks; (iii)~identification of
\textbf{dataset-conditional noise robustness} tied to adaptation frontier
width, with practical guidance on when EWL should and should not be applied;
(iv)~a full ablation over rank, temperature, EMA decay, and class imbalance
(Appendix).
